\addchap{Prefazione}

Queste note nascono come complemento al testo di Ballentine, e non come una trattazione autonoma della meccanica quantistica.

Il Ballentine è il riferimento che seguo per l'impostazione generale della teoria. Su diversi argomenti, però, il libro tace oppure procede per una strada che non è quella vista a lezione: in alcuni casi manca la motivazione fisica da cui si parte in aula, in altri il risultato c'è ma è ottenuto con un metodo differente, in altri ancora l'argomento non compare affatto. Ciò che raccolgo qui è esattamente questo scarto.

Il materiale proviene dai corsi di Meccanica Quantistica tenuti dal professor Ettore Vicari all'Università di Pisa e dal professor Vittorio Giovannetti alla Scuola Normale Superiore.

Da questa impostazione segue la convenzione che adotto ovunque: \emph{di ogni argomento do per acquisito quanto il Ballentine espone fino a quel punto}.
Definizioni, notazioni e dimostrazioni già presenti nel testo non vengono ripetute, e quando servono mi limito a rimandarvi. Di conseguenza questi appunti non si leggono da soli, ma accanto al libro.
